% ============================================================================
%  Chapter 24 -- Beyond One Core                                 (Lecture 24)
% ============================================================================
\chapter{Beyond One Core: Superscalar RISC-V, HPC, GPUs and FPGAs}
\label{ch:beyond}

\epigraphline{If you were plowing a field, which would you rather use: two
strong oxen or 1024 chickens?}{attributed to Seymour Cray}

\begin{chapgoals}
By the end of this chapter you will be able to
\begin{itemize}
\item write the performance equation and say which factor each technique attacks;
\item explain how a superscalar, out-of-order processor executes several
  instructions per cycle, and name \riscv cores that do;
\item read a loop written with the \riscv vector extension;
\item describe a supercomputer, compute a speed-up with Amdahl's law, and say why energy is the limit;
\item say what a GPU and an FPGA are good at, and what each costs the programmer.
\end{itemize}
\end{chapgoals}

This closing chapter is an outlook. It starts from one loop of \vermone and
follows it onto four kinds of machine: the same instructions, on machines
that execute many of them at once. It is not part of the midterm test. The
\riscv vector listings were tested on QEMU; the CUDA and Verilog fragments
are standard textbook forms that were not compiled for this book. The
figures for supercomputers are rounded and come from the TOP500 lists of 2022
to 2025.

% ----------------------------------------------------------------------------
\section{Warm-up on the machine}

\asmfile[caption={The loop that empties the map, from \code{new\_game} of \vermone.},label=lst:by:loop]{code/l24/c03.S}

Listing~\ref{lst:by:loop} sets 960 bytes to zero: four instructions per byte,
3\,840 instructions in all, one byte at a time and one instruction at a time.
At best one instruction per clock cycle and 160 million cycles per second on
the kit: about 24 microseconds. Nobody needs this loop to be faster; but it
is small enough to be followed onto every machine of this chapter. How many
ways are there of making it faster?

\subsection{The performance equation}

\begin{theorem}{Performance equation}
The time that a program takes is the product of three numbers,
\[ T = N\times \mathit{CPI}\times t , \]
where $N$ is the number of instructions executed, $\mathit{CPI}$ the average
number of clock cycles per instruction and $t$ the duration of a clock cycle.
\end{theorem}

For the loop above on the kit: $N=3\,840$, about 1 cycle each, 6.25
nanoseconds per cycle. Each factor can be attacked: \emph{fewer
instructions}, with better algorithms, richer instructions, vectors;
\emph{fewer cycles per instruction}, with pipelines and superscalar execution;
\emph{a shorter cycle}, with faster transistors and deeper pipelines. And
when one processor is not enough: more processors working together. Every
machine of this chapter is a different answer to the same equation
\citep{patterson2020}.

\subsection{Why not simply a faster clock?}

For thirty years each new generation of transistors was smaller, faster, and
no hotter per square millimetre. Round 2005 that stopped: smaller
transistors still arrive, but a faster clock means more heat than a chip can
shed (Chapter~\ref{ch:electronics}), and clock frequencies have stayed at a
few gigahertz ever since. Moore's law, more transistors, went on; the rule
that kept the density of power constant, known as Dennard scaling, ended.
The extra transistors went elsewhere: wider processors, more cores, vector
units, accelerators. Performance now comes from doing many things at once:
from \term{parallelism}. The 160~MHz of the kit are slow by choice: they cost
milliwatts.

% ----------------------------------------------------------------------------
\section{Superscalar processors}

\subsection{From one at a time to many at once}

\begin{table}[htb]
\centering\small
\caption{From one instruction at a time to many at once.}\label{tab:by:kinds}
\begin{tabularx}{\textwidth}{@{}P{0.21\textwidth}LP{0.14\textwidth}L@{}}
\toprule
\textbf{Processor} & \textbf{How it works} & \textbf{Instr.\ per cycle} & \textbf{Example}\\
\midrule
Multi-cycle & one instruction completes before the next begins & less than 1 & the model of Chapter~\ref{ch:processor}\\
Pipelined & several instructions in flight, each in a different stage & at most 1 & the ESP32-C6 of the kit: four stages\\
Superscalar, in order & two or more pipelines side by side & up to 2 or more & small application cores\\
Superscalar, out of order & instructions run as soon as their operands are ready & up to 4, 6, 8, \ldots & laptops, telephones, servers\\
\bottomrule
\end{tabularx}
\end{table}

Chapter~\ref{ch:instructions} introduced the pipeline and named superscalar
processors in one line; this section opens that line
(Table~\ref{tab:by:kinds}). The \emph{width} of a processor is the number of
instructions that it can start in one cycle. With two pipelines side by side
and the five classic stages, six instructions complete in seven cycles
(Table~\ref{tab:by:two}); a single pipeline needs ten, and a processor
without a pipeline thirty.

\begin{table}[htb]
\centering\small
\caption{Two pipelines side by side: six instructions in seven cycles, if
nothing gets in the way.}\label{tab:by:two}
\begin{tabular}{@{}l*{7}{>{\centering\arraybackslash}p{8mm}}@{}}
\toprule
\textbf{cycle} & 1 & 2 & 3 & 4 & 5 & 6 & 7\\
\midrule
$i_1$ & \cellcolor{prgaccent!20}IF & \cellcolor{prgamber!25}ID & \cellcolor{prgviolet!18}EX & \cellcolor{prggreen!20}MEM & \cellcolor{prgred!15}WB &&\\
$i_2$ & \cellcolor{prgaccent!20}IF & \cellcolor{prgamber!25}ID & \cellcolor{prgviolet!18}EX & \cellcolor{prggreen!20}MEM & \cellcolor{prgred!15}WB &&\\
$i_3$ && \cellcolor{prgaccent!20}IF & \cellcolor{prgamber!25}ID & \cellcolor{prgviolet!18}EX & \cellcolor{prggreen!20}MEM & \cellcolor{prgred!15}WB &\\
$i_4$ && \cellcolor{prgaccent!20}IF & \cellcolor{prgamber!25}ID & \cellcolor{prgviolet!18}EX & \cellcolor{prggreen!20}MEM & \cellcolor{prgred!15}WB &\\
$i_5$ &&& \cellcolor{prgaccent!20}IF & \cellcolor{prgamber!25}ID & \cellcolor{prgviolet!18}EX & \cellcolor{prggreen!20}MEM & \cellcolor{prgred!15}WB\\
$i_6$ &&& \cellcolor{prgaccent!20}IF & \cellcolor{prgamber!25}ID & \cellcolor{prgviolet!18}EX & \cellcolor{prggreen!20}MEM & \cellcolor{prgred!15}WB\\
\bottomrule
\end{tabular}
\end{table}

\subsection{What gets in the way: dependences}

\asmfile[caption={Dependences between neighbouring instructions.},label=lst:by:dep]{code/l24/c04.S}

This is the ideal case, and it is rare. In Listing~\ref{lst:by:dep} the
instruction $i_2$ cannot run beside $i_1$, because it uses its result: a
\term{data dependence}. $i_3$ depends on nobody: why should it wait? And
after the branch $i_4$ the processor does not know what comes next: a
\term{control dependence}. Real code is full of such chains, and width alone
is not enough. A superscalar processor that works in order stalls at $i_2$,
and everything behind it waits. The remedy is to let $i_3$ overtake.

\subsection{Out-of-order execution}

\begin{enumerate}
\item \textbf{Fetch and decode, in order.} Several instructions per cycle enter
  a window of tens or hundreds of instructions.
\item \textbf{Rename.} Each result receives a fresh physical register, so that
  two uses of the name \reg{t0} no longer collide.
\item \textbf{Execute when ready.} An instruction starts as soon as its
  operands exist, whatever its position in the program.
\item \textbf{Commit, in order.} Results become official in the order of the
  program: seen from outside, nothing ever overtook.
\end{enumerate}

The algorithm is Tomasulo's, designed in 1967 for the floating-point unit of
the IBM System/360 Model~91; the reorder buffer that commits in order came in
the 1980s. The programmer's model, the one learnt in this book, is preserved
exactly: this is the contract of the instruction set.

\subsection{Register renaming: thirty-two names, many registers}

\asmfile[caption={The same code as written, and as executed after renaming.},label=lst:by:rename]{code/l24/c05.S}

In Listing~\ref{lst:by:rename} the second half reuses the \emph{name}
\reg{t0}, not its value: a \term{false dependence}. Once renamed, the two
halves share nothing, and can run at the same time on different registers.
(The numbers \code{p41} to \code{p44} are invented.) The thirty-two registers
known to the programmer are names; the chip has many more, a few hundred in a
large core.

\subsection{Guessing the future: branch prediction}

A wide processor meets a branch every few instructions and cannot wait to
know its outcome. It \emph{predicts} the direction from the history of that
branch, and goes on executing along the guess: a right guess costs nothing; a
wrong one throws away the speculative work and starts again. The loop of the
warm-up branches back 959 times and falls through once: a predictor learns it
in the first turns. Modern predictors are right far more than nine times out
of ten.

\begin{prgaside}[Speculation and security]
Speculative execution has a dark side. In 2018 the attacks called Spectre
and Meltdown showed that work which is thrown away can still leave traces in
the cache that reveal secrets. Tricks of performance are also questions of
security.
\end{prgaside}

\subsection{RISC-V goes superscalar}

\begin{table}[htb]
\centering\small
\caption{Some \riscv cores. The list ages quickly.}\label{tab:by:cores}
\begin{tabularx}{\textwidth}{@{}lP{0.3\textwidth}L@{}}
\toprule
\textbf{Core} & \textbf{From} & \textbf{Kind}\\
\midrule
HP core of the ESP32-C6 & Espressif & in order, four stages, one instruction per cycle: the kit\\
BOOM & University of California, Berkeley & out of order, open source, for research and teaching\\
XiangShan & Chinese Academy of Sciences & out of order, open source, high performance\\
XuanTie C910 & T-Head (Alibaba) & out of order, three-wide; released as open source\\
Performance P550 & SiFive & out of order, three-wide, thirteen stages; sold as a design\\
Veyron & Ventana Micro Systems & out of order, for data centres\\
\bottomrule
\end{tabularx}
\end{table}

Because the instruction set is open, anyone may build such a core
(Table~\ref{tab:by:cores}), and several are published as source code that a
student can read: this has no precedent among commercial architectures.

\begin{definition}{Architecture and microarchitecture}
The \term{architecture} is the contract: the instructions, the registers,
what each instruction does. The \term{microarchitecture} is one way of
honouring it: stages, width, predictors, caches.
\end{definition}

The \ins{add} written for a microcontroller at 160~MHz is the \ins{add} of a
server core, and \vermone would run unchanged on any of the cores of the
table, given a screen. This is why an instruction set learnt once is
knowledge that lasts. The distinction organises the whole book: Chapters
\ref{ch:firstcontact} to~\ref{ch:verm1p4} taught the architecture, and
the kit is one small microarchitecture.

% ----------------------------------------------------------------------------
\section{Vectors: one instruction, many data}

\subsection{Flynn's taxonomy}

\begin{table}[htb]
\centering\small
\caption{The taxonomy of \citet{flynn1972}.}\label{tab:by:flynn}
\begin{tabular}{@{}llll@{}}
\toprule
\textbf{Class} & \textbf{Instructions} & \textbf{Data} & \textbf{Example}\\
\midrule
SISD & one stream & one stream & the kit: one core, one datum at a time\\
SIMD & one stream & many at once & vector instructions; a GPU, to a first approximation\\
MISD & many & one stream & rare: redundant systems that vote\\
MIMD & many & many & multicore processors, clusters, supercomputers\\
\bottomrule
\end{tabular}
\end{table}

The classification of Table~\ref{tab:by:flynn} is crude, and is still the
first vocabulary of parallel computing.

\subsection{The same loop, with vectors}

\asmfile[caption={Emptying the map with the vector extension.},label=lst:by:vec]{code/l24/c06.S}

The V extension of \riscv, ratified in 2021, adds \term{vector registers},
many bytes wide (Listing~\ref{lst:by:vec}). The instruction \ins{vsetvli}
asks: how many of the bytes that are left fit in a register? The answer
comes back in \reg{t2}. Then \ins{vse8.v} stores \reg{t2} bytes with one
instruction. With registers of 128 bits the loop handles 16 bytes per turn:
60 turns instead of 960.

The code never says how wide the registers are: it is \emph{vector-length
agnostic}, and the same binary uses 128, 256 or 1\,024 bits. The extension is
not on the kit, since the ESP32-C6 is RV32IMAC; the listing was tested on
QEMU with registers of 128, 256, 512 and 1\,024 bits.

\subsection{A sum, lane by lane}

\asmfile[caption={The sum of sixteen words with vectors.},label=lst:by:vsum]{code/l24/c07.S}

Listing~\ref{lst:by:vsum} is the parallel sum of Chapters
\ref{ch:information} and~\ref{ch:video}, in \riscv. \ins{vle32.v} loads
several words at once; \ins{vadd.vv} adds register to register, \emph{lane}
by lane; at the end \ins{vredsum.vs} adds the lanes together. With the
sixteen numbers from 1 to 16 the result is 136. (The setting \code{tu},
\emph{tail undisturbed}, keeps the partial sums of the lanes that the last
turn does not use.)

% ----------------------------------------------------------------------------
\section{High-performance computing}

\subsection{What is a supercomputer?}

A supercomputer is not one enormous processor. It is thousands of ordinary
computers, the \term{nodes}, working on the same problem
(Figure~\ref{fig:by:cluster}). Each node has one or two processors with tens
of cores, often several GPUs, and a great deal of memory. A very fast network
joins the nodes, and a parallel file system serves them all. Speed is
measured in floating-point operations per second, FLOPS, on a standard
problem: a huge system of linear equations (the benchmark called HPL,
\emph{High-Performance Linpack}). The TOP500 list has ranked the fastest
machines twice a year since 1993.

\begin{figure}[htb]
\centering
\includegraphics[width=0.8\textwidth]{rcf_cluster.jpg}
\caption{Racks of nodes of the Research Computing Facilities of the University of
Naples ``Parthenope'': the structure of a supercomputer.}
\label{fig:by:cluster}
\end{figure}

\begin{table}[htb]
\centering\small
\caption{Orders of magnitude (rounded HPL figures from the TOP500 lists; the
figure for the laptop is an order of magnitude).}\label{tab:by:flops}
\begin{tabularx}{\textwidth}{@{}llL@{}}
\toprule
\textbf{Machine} & \textbf{Where} & \textbf{Floating-point operations per second}\\
\midrule
ESP32-C6 & the kit & no floating-point unit: fixed point, as in \vermone\\
A laptop & your bag & about $10^{11}$\\
Leonardo (2022) & CINECA, Bologna & about $2.4\times 10^{17}$\\
Frontier (2022) & Oak Ridge, USA & over $10^{18}$: the first exascale machine\\
El Capitan (2024) & Livermore, USA & about $1.7\times 10^{18}$\\
JUPITER (2025) & J\"ulich, Germany & about $10^{18}$: the first European exascale machine\\
\bottomrule
\end{tabularx}
\end{table}

\emph{Exa} is $10^{18}$. Seven orders of magnitude separate a laptop from an
exascale machine (Table~\ref{tab:by:flops}): about fifty years of Moore's
law, bought with parallelism instead.

\subsection{Amdahl's law}

\begin{theorem}{Amdahl's law}
If a fraction $p$ of the work of a program can be done in parallel, the
speed-up obtained with $n$ processors is
\[ S(n)=\frac{1}{(1-p)+\dfrac{p}{n}},\qquad\text{and}\qquad
   \lim_{n\to\infty}S(n)=\frac{1}{1-p}. \]
\end{theorem}

\noindent\emph{Proof.} Take the time on one processor as 1. The serial part
takes $1-p$ whatever the number of processors; the parallel part, shared
among $n$, takes $p/n$. The speed-up is the ratio of the two times.
\hfill$\square$

\smallskip
The serial fraction does not shrink, however many processors are added
\citep{amdahl1967}. With $p=0.9$ the speed-up can never exceed 10, even with
a million processors (Table~\ref{tab:by:amdahl}): a thousand processors on a
program that is 90\% parallel make it less than ten times faster. Before
buying processors, measure $p$. The same law holds for any improvement that
applies to only a part of the work: a faster GPU does not help the part that
runs on the CPU.

\begin{table}[htb]
\centering\small
\caption{Amdahl's law in numbers: the speed-up.}\label{tab:by:amdahl}
\begin{tabular}{@{}ccccc@{}}
\toprule
\textbf{Parallel fraction $p$} & $n=10$ & $n=100$ & $n=1000$ & $n$ unlimited\\
\midrule
0.50 & 1.8 & 2.0 & 2.0 & 2\\
0.90 & 5.3 & 9.2 & 9.9 & 10\\
0.99 & 9.2 & 50 & 91 & 100\\
0.999 & 9.9 & 91 & 500 & 1000\\
\bottomrule
\end{tabular}
\end{table}

As a check, $p=0.9$ and $n=100$ give $1/(0.1+0.009)=9.17$. Gustafson's
answer, in 1988, was that with more processors people solve bigger problems,
and that in a bigger problem the parallel part usually grows faster than the
serial one.

\subsection{Two ways to share the work}

With \term{shared memory}, several cores of one node see the same memory, and
the program runs as \emph{threads} that read and write the same variables.
It is easy to start: OpenMP adds one line above a loop. The danger is two
threads that write the same variable. With \term{distributed memory} each
node has its own memory and nothing is shared; the program runs as
\emph{processes} that send each other messages, and MPI, the \emph{Message
Passing Interface}, is the standard. It scales to a million cores, and the
programmer pays for it. Real applications use both: MPI between the nodes,
threads and GPUs inside each node. The weather and marine models of the
Department work in this way.

\subsection{The limit is energy}

An exascale machine draws about twenty megawatts, the consumption of a small
town (Frontier reached its exaflop at about 21~MW): some 50 thousand million
operations per joule. The electricity of its lifetime costs as much as the
machine. So the figure that matters is operations per watt, which has its
own ranking, the Green500, led by GPUs and other accelerators: specialised
hardware does more with each joule. The same reason puts a small \riscv core
in the kit: milliwatts. Energy efficiency is the thread that ties a
microcontroller to a supercomputer.

\subsection{RISC-V and European supercomputing}

Today's largest machines use x86 or Arm processors with GPU accelerators.
Europe designs its own processors and accelerators on \riscv: the European
Processor Initiative, and DARE, the project named in
Chapter~\ref{ch:intro} \citep{eurohpcdare}. Monte Cimone, built in 2022 at
the University of Bologna with CINECA, is a complete HPC cluster on \riscv
processors; the Barcelona Supercomputing Center designs and fabricates
\riscv cores of its own. An open instruction set lets a continent build what
it would otherwise have to buy.

% ----------------------------------------------------------------------------
\section{GPU accelerators}

\subsection{Two ways to spend transistors}

A CPU has a few large cores, tens of them, each out of order, with predictors
and big caches: it is built to finish \emph{one} task as soon as possible,
and is at ease with branches and irregular data. A GPU has thousands of
small, simple cores, all executing the same code on different data: it is
built to finish a million tasks per second, and is at ease with regular
arrays: pixels, matrices, grids. These are Cray's oxen and chickens, and the
answer depends on the field. The GPU is not faster at any single thing; it
does more things at once.

\subsection{How a program uses a GPU}

\begin{enumerate}
\item \textbf{Prepare.} The program on the CPU, the \emph{host}, builds its
  arrays in main memory.
\item \textbf{Copy in.} The data travel to the memory of the GPU: a real cost,
  in time.
\item \textbf{Launch.} One small function, the \term{kernel}, runs in
  thousands of threads, one per element of the data.
\item \textbf{Copy out.} The results travel back, and the host goes on.
\end{enumerate}

The GPU is a device, like those of Chapter~\ref{ch:io}: the CPU sends it data
and commands through a driver. The programming models are CUDA, OpenCL and
SYCL, and directives such as OpenACC and OpenMP offload.

\cfile[caption={The map of \vermone with one thread per cell (CUDA C, in its
usual textbook form; not compiled for this book).},label=lst:by:cuda]{code/l24/c08.c}

In Listing~\ref{lst:by:cuda} there is no loop: there are 960 threads. Each
thread computes its own index and does one store. Threads come in blocks,
and the last ones find nothing to do. For 960 bytes the two copies cost far
more than the work, and the point is serious: a GPU pays off when the data
are large and the work per element is real.

\subsection{When does a GPU pay off?}

It pays off when the same operation applies to very many data, independently:
images, matrices, the grids of a simulation, neural networks; and when the
data can stay on the GPU for many steps instead of travelling at every one.
It does not when the work is a chain of decisions, each depending on the one
before. And Amdahl's law applies again: what stays on the CPU sets the limit.
A weather simulation updates every cell of a three-dimensional grid from its
neighbours, millions of times: the same shape as the map of \vermone, a
little larger. The training of today's neural networks is, almost entirely,
matrix arithmetic on GPUs.

% ----------------------------------------------------------------------------
\section{FPGA accelerators}

\subsection{A chip that becomes the circuit you describe}

\begin{definition}{FPGA}
A \term{field-programmable gate array} (FPGA) is a chip made of a sea of small
programmable blocks: \term{look-up tables}, each able to be any function of a
few inputs; flip-flops, one beside each table; and wires and switches that can
join any block to any other. A configuration file, the \term{bitstream}, sets
every table and every switch.
\end{definition}

Load another bitstream, and the same chip is another circuit. The first FPGA,
the Xilinx XC2064, was sold in 1985; today they go from a few euros to boards
for data centres. An FPGA does not execute a program: it \emph{becomes} a
circuit. Nothing is fetched or decoded, and all the gates work at the same
time.

\subsection{A look-up table is a truth table in memory}

\begin{center}\small\ttfamily
\begin{tabular}{cc|ccc}
\textrm{$a$}&\textrm{$b$}&\textrm{stored for AND}&\textrm{stored for XOR}&\textrm{stored for NOT $a$}\\\hline
0&0&0&0&1\\0&1&0&1&1\\1&0&0&1&0\\1&1&1&0&0
\end{tabular}
\end{center}

The inputs are the address, and four stored bits choose the function:
changing the bits changes the gate. Chapter~\ref{ch:boole1} said that a
truth table defines a function, and Chapter~\ref{ch:storage} that a ROM
addressed by the inputs implements any function. A look-up table is that
ROM, with four to six inputs and sixteen to sixty-four bits; the flip-flop
next to it is that of Chapter~\ref{ch:sequential}. An FPGA is this book,
repeated a hundred thousand times on one chip.

\subsection{The U-turn rule of Verm1, as hardware}

\lstinputlisting[style=verilog,caption={The U-turn rule in Verilog (not synthesised for this book).},label=lst:by:verilog]{code/l24/c09.v}

Verilog, like VHDL, is a \term{hardware description language}: it describes
circuits, not steps (Listing~\ref{lst:by:verilog}). The statement
\code{assign} is a wire, always true: there is no before and no after. The
circuit is two XOR gates and an AND with one inverted input,
\[ \mathit{refuse}=(d_1\oplus w_1)\cdot\nott{(d_0\oplus w_0)} . \]
On the kit this rule took three instructions and three cycles
(Chapter~\ref{ch:verm1p2}); here the gates answer in a fraction of a
nanosecond, continuously.

\subsection{From description to circuit}

\begin{figure}[htb]
\centering
\begin{tikzpicture}[node distance=4mm]
  \node[blkacc,text width=21mm,minimum height=16mm] (a) {\textbf{Describe}\\[1pt]\footnotesize Verilog, VHDL, or C};
  \node[blkamb,text width=21mm,minimum height=16mm,right=of a] (b) {\textbf{Simulate}\\[1pt]\footnotesize waveforms};
  \node[blkvio,text width=21mm,minimum height=16mm,right=of b] (c) {\textbf{Synthesise}\\[1pt]\footnotesize tables and flip-flops};
  \node[blkgrn,text width=21mm,minimum height=16mm,right=of c] (d) {\textbf{Place and route}\\[1pt]\footnotesize on the chip};
  \node[blkred,text width=21mm,minimum height=16mm,right=of d] (e) {\textbf{Configure}\\[1pt]\footnotesize load the bitstream};
  \foreach \x/\y in {a/b,b/c,c/d,d/e} \draw[->,thick] (\x) -- (\y);
\end{tikzpicture}
\caption{From a description to a circuit.}
\label{fig:by:fpgaflow}
\end{figure}

Each attempt takes minutes to hours (Figure~\ref{fig:by:fpgaflow}): slower
than building a cartridge, far faster than making a chip. The same flow,
with a factory at the end instead of a bitstream, produces an \term{ASIC}, a
chip made for one purpose. FPGAs are where chips are tried before they are
made.

\subsection{A processor is a circuit too}

A \riscv core is a description in a hardware language, and can be loaded on
an FPGA. Such \term{soft cores} are free to read and to use: PicoRV32,
VexRiscv, NEORV32, Rocket. A core of the class of the kit's fits in a small,
cheap FPGA, and beside it one can put a circuit of one's own for the one
operation that a program repeats most: a custom accelerator. \riscv reserves
room in its encoding for custom instructions, to reach such circuits.

\begin{prgconcept}[The circle closes]
The gates of Chapter~\ref{ch:boole1} build the adder of
Chapter~\ref{ch:boole2}, the registers of Chapter~\ref{ch:registers}
and the processor of Chapter~\ref{ch:processor}; and that processor is a file
that can be read. From gates to a processor, and back: the reader now knows
every layer in between.
\end{prgconcept}

\subsection{Four ways to compute}

\begin{table}[htb]
\centering\small
\caption{Four ways to compute.}\label{tab:by:four}
\begin{tabularx}{\textwidth}{@{}lLLLL@{}}
\toprule
 & \textbf{CPU} & \textbf{GPU} & \textbf{FPGA} & \textbf{ASIC}\\
\midrule
What it is & a few complex cores & thousands of simple cores & a circuit you configure & a circuit made for one job\\
Programmed in & any language & kernels: CUDA, OpenCL, SYCL & Verilog, VHDL, HLS & the same, then a factory\\
Strong at & everything, moderately & regular data in bulk & streams, low latency & one thing, perfectly\\
Work per joule & low & high & high & highest\\
Effort to use & low & medium & high & very high\\
Change of mind & recompile & recompile & a new bitstream & a new chip\\
\bottomrule
\end{tabularx}
\end{table}

Generality is paid in energy; efficiency is paid in effort
(Table~\ref{tab:by:four}). A modern system on a chip has all four: a
telephone holds CPU cores, a GPU, and fixed circuits for video and radio.

% ----------------------------------------------------------------------------
\section{Back to the machine}

\begin{table}[htb]
\centering\small
\caption{Emptying the map of \vermone, 960 cells, on five machines.}\label{tab:by:five}
\begin{tabularx}{\textwidth}{@{}P{0.27\textwidth}LP{0.22\textwidth}@{}}
\toprule
\textbf{Machine} & \textbf{How} & \textbf{Steps in time}\\
\midrule
The kit: ESP32-C6 & one byte per turn, one instruction per cycle & 3\,840 cycles\\
Out-of-order \riscv core & the same binary; several instructions per cycle, the branch predicted & a fraction of those cycles\\
\riscv with vectors, 128 bits & 16 bytes per turn & 60 turns\\
GPU & 960 threads, one per cell & one launch, plus two copies\\
FPGA & 960 flip-flops with a common reset & one clock cycle\\
\bottomrule
\end{tabularx}
\end{table}

The first three rows of Table~\ref{tab:by:five} are the same instruction set;
the first and the third were run for this book. The work is the same
everywhere; what changes is how much of it happens at the same time. And the
costs change: the first row needs a chip of two euros and four lines of
assembly; the last needs a hardware design.

\begin{prgconcept}
The engineer's question is never which machine is fastest, but which is
enough.
\end{prgconcept}

% ----------------------------------------------------------------------------
% Generated by tools/textbook/export_listings.js from decks/plan.js and code/l24.S: do not edit.
\section{The program of this chapter}\label{sec:lab:l24}

The listings of this chapter are parts of one program, \fname{code/l24.S} in the
repository of the course, printed in full in Listing~\ref{lst:lab:l24}. The
automatic tests of the course material build it and run it on the headless
runner of \prgqt{}, with scripted input, and check the frame;
Figure~\ref{fig:lab:l24} is the frame that the test wrote.

\begin{description}[style=nextline,leftmargin=1.5em]
\item[On the screen.] A green square: the 960 bytes of the map, filled with ones by init, were all set to zero by the four-instruction loop of the lecture. A red square would mean that some byte survived.
\item[What to press.] Nothing: the program runs by itself.
\end{description}

\begin{figure}[htb]
\centering
\includegraphics[width=0.5\textwidth]{lab_l24.png}
\caption{\fname{l24.S} on \prgqt{}.}\label{fig:lab:l24}
\end{figure}

\asmfile[caption={The complete program of the chapter, \fname{code/l24.S}.},label=lst:lab:l24]{code/l24/program.S}

To run it on \prgqt{} the commands are those of Listing~\ref{lst:lab:l24:run};
Section~\ref{sec:prg:testing} gives the same steps for QEMU and for the board.

\begin{shell}[caption={Building \fname{l24.S} and running it on \prgqt{}.},label=lst:lab:l24:run]
# 1. start PRG32-QT on this computer (it listens on port 8080)

# 2. in the PRG32 folder, with ESP-IDF active: build the cartridge
python3 -m prg32 cartridge build ~/arc12/code/l24.S --portable \
    --entry-prefix l24 --name l24 --out l24.prg32

# 3. send it to PRG32-QT and start it
python3 -m prg32 esp32c6 upload-and-run l24.prg32 \
    --url http://localhost:8080 --slot cart0

# 4. look, press, compare with the previous slide
\end{shell}


% ----------------------------------------------------------------------------
\section{Check your understanding}

\begin{quiz}{A program is 80\% parallel. What is the largest speed-up that any number of processors can give?}
\item 4
\item 5
\item 8
\item 80
\end{quiz}

\begin{quiz}{What does register renaming remove?}
\item branches
\item the need for a stack
\item false dependences between uses of the same register name
\item cache misses
\end{quiz}

\begin{quiz}{With 256-bit vector registers, how many turns does the vector loop need for the 960 bytes?}
\item 4
\item 30
\item 60
\item 960
\end{quiz}

\begin{quiz}{A two-input look-up table stores, for the inputs 00, 01, 10, 11, the bits 0, 1, 1, 1. Which gate is it?}
\item AND
\item OR
\item XOR
\item NAND
\end{quiz}

\section{Exercises}

\begin{exercise}[paper first, then the machine]
\begin{enumerate}[nosep]
\item Fill in the Amdahl table for $p=0.95$ and $n=10$, 100, 1000.
\item Draw the pipeline diagram of a three-wide processor for nine independent
  instructions with five stages.
\item Rename the registers of the code of \code{step} that moves and wraps the
  head in \vermone, Part~2: which instructions could run in the same cycle?
\item Write the truth table stored in a four-input look-up table for
  $\mathit{refuse}=(\mathit{dir}\oplus\mathit{want})=2$.
\item Build the test \fname{vector\_qemu.S} of the course material and run it
  on QEMU with vector registers of 128 and of 512 bits: why does the same
  binary work?
\end{enumerate}
\expected{Amdahl with $p=0.95$: 6.9, 16.8, 19.6. Nine instructions on a
three-wide processor: seven cycles. The look-up table has four ones out of
sixteen entries, for $(\mathit{dir},\mathit{want})=(0,2)$, $(1,3)$, $(2,0)$,
$(3,1)$. QEMU is started as
\code{qemu-system-riscv32 -machine virt -cpu rv32,v=true,vlen=128}.}
\end{exercise}

\section{Glossary}

\begin{glossarytable}
superscalar & superscalare\\
register renaming & ridenominazione dei registri\\
out-of-order execution & esecuzione fuori ordine\\
branch prediction & predizione dei salti\\
speed-up & accelerazione (speed-up)\\
node & nodo\\
shared memory & memoria condivisa\\
message passing & scambio di messaggi\\
accelerator & acceleratore\\
kernel (GPU) & kernel\\
look-up table & tabella di look-up\\
bitstream & bitstream\\
hardware description language & linguaggio di descrizione dell'hardware\\
soft core & soft core\\
\end{glossarytable}

\section{Conclusions and reading}

One equation, three factors: every fast machine attacks one of them, or
multiplies the processors. Superscalar, out-of-order cores keep the contract
of the instruction set and execute many instructions at once. Vectors, GPUs
and clusters exploit data that can be processed independently, and Amdahl's
law says how far that goes. FPGAs turn a description into a circuit: the
logic of this book, at scale.

The book ends here: from a bit to a supercomputer, with one instruction set.

\paragraph{Reading.} \citet{patterson2020}, chapters 4 and~6;
\citet{hennessy2019}; \url{https://top500.org}; the \riscv vector
specification at \url{https://riscv.org}.
